\subsection{面片}

不属于超平面集 H 中任何超平面的 E 中点集是开集，因为 H 是局部有限的。更精确地，有如下结果：

命题 1. 设 a 是 E 的一点。存在 a 的一个连通开邻域，它不与属于 H 且不经过 a 的任何超平面相交。此外，属于 H 且经过 a 的超平面只有有限多个。

满足 \(H \in H\) 且 \(a \notin H\) 的超平面 H 所成的集合 N 是局部有限的，因为它包含于 H 中。因此，不属于 N 中任何超平面的 E 中点集 U 是开集。由于 \(a \in U\)，存在包含于 U 的 a 的连通开邻域。命题的其余部分显然。

给定两点 \(x\) 和 \(y\)，它们属于 \(\mathbf{E}\)，用 \(\mathrm{R}\{x,y\}\) 表示关系

“对于任意超平面 \(H \in H\)，或者 \(x \in H\) 且 \(y \in H\)，或者 x 和 y 严格位于 H 的同侧。”

显然，R 是 E 上的等价关系。

定义 1. 相对于 \(\mathfrak{H}\) 的 E 的一个面片，是上述等价关系的一个等价类。

命题 2. 面片集是局部有限的。

这显然是因为 \(\mathfrak{H}\) 是局部有限的。

设 F 是一个面片，\(a\) 是 F 的一点。超平面 \(H \in H\) 包含 F 当且仅当 \(a \in H\)；因此这些超平面所成的集合 F 是有限的；它们的交是 E 的一个仿射子空间 L，我们称其为 F 的仿射支撑；L 的维数称为 F 的维数。

若 \(\mathfrak{N}\) 是超平面 \(H\in \mathfrak{H}\) 中不包含 \(F\) 的那些超平面所成的集合，则

\[
F = L \cap \bigcap_ {H \in \mathfrak {N}} D _ {H} (a).\tag{2}
\]

我们将证明 F 的闭包由下式给出：

\[
\overline {{F}} = L \cap \bigcap_ {H \in \mathfrak {N}} \overline {{D _ {H} (a)}}.\tag{3}
\]

显然，右端包含左端。反之，设 \(x \in L \cap \bigcap_{H \in \mathfrak{N}} \overline{\mathrm{D}_H(a)}\)。端点为 \(a\) 和 \(x\) 的开线段包含于 \(L\) 以及每个 \(D_H(a)\)（对于 \(H \in \mathfrak{N}\)）中，因而包含于 \(F\) 中。于是 \(x\) 属于 \(F\) 的闭包，公式得证。

命题 3. 设 F 是一个面片，L 是其仿射支撑。

\begin{enumerate}
\def\labelenumi{(\roman{enumi})}
\item
  集合 \(\mathbf{F}\) 是 \(\mathbf{L}\) 的仿射子空间 \(\mathbf{E}\) 的凸开子集。
\item
  \(\mathbf{F}\) 的闭包是 \(\mathbf{F}\) 与维数严格小于 \(\mathbf{F}\) 的面片之并。
\item
  在拓扑空间 L 中，集合 F 是其闭包的内部。
\end{enumerate}

由于每个开半空间和每个超平面都是 E 的凸子集，公式 (2) 表明 F 是一族凸子集的交，因而是凸的。另一方面，设 a 属于 F，取 E 中 a 的一个凸开邻域 U，使其不与集合 N 中的任何超平面相交，其中 N 由满足 \(H \in H\) 且 \(a \notin H\) 的 H 组成。于是对于任意 \(H \in N\)，有 \(U \subset D_{H}(a)\)，从而 \(L \cap U \subset F\)，所以 F 在拓扑空间 L 中是开的。

设 \(b\) 是 \(\overline{F} - F\) 中属于一个面片 \(F'\) 的点，令 \(\mathfrak{N}'\) 为经过 \(H \in \mathfrak{N}\) 的、属于 \(b\) 的超平面所成的集合；置 \(\mathfrak{N}'' = \mathfrak{N} - \mathfrak{N}'\)。对于 \(H\) 中的任意 \(\mathfrak{N}''\)，有 \(b \notin H\) 且 \(b \in \overline{D_H(a)}\)，因而 \(b \in D_H(a)\) 且 \(D_H(b) = D_H(a)\)；由面片的定义，故有

\[
F ^ {\prime} = L \cap \bigcap_ {H \in \mathfrak {N} ^ {\prime}} H \cap \bigcap_ {H \in \mathfrak {N} ^ {\prime \prime}} D _ {H} (a),\tag{4}
\]

而 (3) 蕴含

\[
\overline {{\mathbf {F}}} = \mathbf {L} \cap \bigcap_ {\mathsf {H} \in \mathfrak {N} ^ {\prime}} \overline {{\mathsf {D} _ {\mathsf {H}} (a)}} \cap \bigcap_ {\mathsf {H} \in \mathfrak {N} ^ {\prime \prime}} \overline {{\mathsf {D} _ {\mathsf {H}} (a)}},\tag{5}
\]

于是 \(\mathbf{F}' \subset \overline{\mathbf{F}}\)。不能有 \(\mathfrak{N}' = \varnothing\)，因为由 (2) 和 (4) 这将推出 \(\mathbf{F} = \mathbf{F}'\)，这与假设 \(b \notin \mathbf{F}\) 且 \(b \in \mathbf{F}'\) 矛盾。\(\mathbf{F}'\) 的支撑是集合 \(\mathbf{L}' = \mathbf{L} \cap \bigcap_{\mathbf{H} \in \mathfrak{N}'} \mathbf{H}\)；我们有 \(a \in \mathbf{L}\)，但对于 \(a \notin \mathbf{H}\) 中的 \(\mathbf{H}\)，有 \(\mathfrak{N}'\)，所以 \(\mathbf{L}' \neq \mathbf{L}\)，最终 \(\dim \mathbf{L}' < \dim \mathbf{L}\)，即 \(\dim \mathbf{F}' < \dim \mathbf{F}\)。这证明了 (ii)。

设 H 属于 \(N'\)，令 D 为由 H 界定且不同于 \(D_{\mathrm{H}}(a)\) 的开半空间；有 \(b \in H \cap L\)，并且立即可见，\(D \cap L\) 是 L 中由 L 的超平面 \(H \cap L\) 界定的半空间；因此，L 中 b 的每个邻域都与 \(D \cap L\) 相交，而由 (3)，\(D \cap L\) 与 \(\overline{F}\) 不相交，于是可见 \(\overline{F} - F\) 中的点 b 不可能是拓扑空间 L 中 \(\overline{F}\) 的内部点。由于 F 在 L 中开，得到 (iii)。证毕。

推论. 设 F 和 \(F'\) 是两个面片。若 \(\overline{F} = \overline{F'}\)，则面片 F 和 \(F'\) 相等。

这由 (iii) 得出。

命题 4. 设 F 是一个面片，L 是 E 的一个仿射子空间，且它是属于 H 的超平面的交：记 N 为不包含 L 的超平面 \(H \in H\) 所成的集合。下列条件等价：

\begin{enumerate}
\def\labelenumi{(\roman{enumi})}
\item
存在面片 \(\mathbf{F}'\)，它的支撑为 L 且与 \(\overline{\mathbf{F}}\) 相交。
\item
存在面片 \(\mathbf{F}'\)，其支撑为 \(\mathbf{L}\) 且包含于 \(\overline{\mathbf{F}}\)。
\item
存在点 \(x\) 属于 \(\mathbf{L} \cap \overline{\mathbf{F}}\)，但不属于 \(\mathfrak{N}\) 中的任何超平面。
\end{enumerate}

若这些条件满足，\(\mathrm{L} \cap \mathrm{D}_{\mathfrak{M}}(\mathrm{F})\) 是支撑为 \(\mathrm{L}\) 且包含于 \(\overline{\mathbf{F}}\) 的唯一面片。

\begin{enumerate}
\def\labelenumi{(\roman{enumi})}
\item
\(\Longrightarrow\) (ii)：由于 \(\overline{\mathbf{F}}\) 是面片的并（命题 3，(ii)），每个与 \(\overline{\mathbf{F}}\) 相交的面片都与包含于 \(\overline{\mathbf{F}}\) 的一个面片相交，因而与该面片相等。
\item
\(\Longrightarrow\) (iii)：对每个 \(x\) 属于 \(F'\)，(iii) 成立，因为 \(\mathfrak{H}\) 中包含 \(x\) 的每个超平面都包含 \(F'\)，因而也包含 \(L\)。
\item
\(\Longrightarrow\) (i)：设 \(x\) 是满足 (iii) 的一点，并令 \(\mathbf{F}'\) 为包含 \(x\) 的面片；显然，\(\mathbf{F}'\) 与 \(\overline{\mathbf{F}}\) 相交。令 \(\mathrm{H} \in \mathfrak{H}\)；则 \(x \notin \mathrm{H}\) 当 \(\mathrm{H} \in \mathfrak{N}\) 时成立，而显然 \(x \in \mathrm{H}\) 当 \(\mathrm{H} \notin \mathfrak{N}\) 时成立；因此，\(\mathbf{F}'\) 的支撑是 \(\mathfrak{H} - \mathfrak{N}\) 中超平面的交，并且等于 \(\mathrm{L}\)。
\end{enumerate}

最后，设 \(F'\) 是一个支撑为 L 且包含于 \(\overline{F}\) 的面片，x 是 \(F'\) 的一点；由于没有 \(N \subset H\) 中的超平面经过 x，命题 1 表明存在 x 的一个凸开邻域 U，使其不与 N 中任何超平面相交。由于 x 属于 F 的闭包，\(U \cap F \neq \varnothing\)；现在 N 是不包含 \(H \in H\) 的超平面 \(F'\) 所成的集合，对于 N 中的所有 H，有 \(D_{H}(x) = D_{H}(U) = D_{H}(U \cap F) = D_{H}(F)\)，由公式 (2) 得

\[
\mathrm{F} ^ {\prime} = \mathrm{L} \cap \mathrm{D} _ {\mathfrak {N}} (\mathrm{F}). \quad \text { Q.E.D. }
\]

